% Generated draft: requires mathematical and source review.
\documentclass[11pt]{article}
\usepackage[margin=25mm]{geometry}
\usepackage{fontspec}
\setmainfont{Latin Modern Roman}
\usepackage{amsmath,amssymb,mathtools,unicode-math}
\setmathfont{Latin Modern Math}
\usepackage{enumitem,needspace,xurl,hyperref}
\setlength{\parindent}{0pt}
\setlength{\parskip}{5pt}
\setlength{\emergencystretch}{3em}
\newcommand{\N}{\mathbb{N}}
\newcommand{\Z}{\mathbb{Z}}
\newcommand{\Q}{\mathbb{Q}}
\newcommand{\R}{\mathbb{R}}
\newcommand{\C}{\mathbb{C}}
\newcommand{\F}{\mathbb{F}}
\newcommand{\A}{\mathbb{A}}
\newcommand{\PP}{\mathbb{P}}
\newcommand{\E}{\mathbb{E}}
\newcommand{\eps}{\varepsilon}
\newcommand{\dd}{\,\mathrm{d}}
\begin{document}
\section*{UnsolvedMath ID: 30005996}
\textbf{Source problem code:} catalogue entry
\subsection{Definitions and mathematical statement}
The supplied record asks about an extremal solution \(u^\star\) near an isolated singularity located at the origin. It does not specify the semilinear elliptic equation, the admissible nonlinearities and domains, the boundary conditions, or the solution family whose extremal member is intended. Consequently, the following makes the requested estimate precise as a conditional formulation; it does not supply the missing PDE hypotheses.

For purposes of interpreting the displayed expression, let \(n\geq 1\) denote the unspecified ambient dimension, let \(\Omega\subset\R^n\) be an open neighborhood of \(0\), and write \(|x|\) for the Euclidean norm. Let \(I\subset\R\) be an interval on which \(f:I\to\R\) is differentiable, and let \(u^\star:\Omega\setminus\{0\}\to I\) denote a pointwise representative of the intended extremal solution. These are interpretive requirements for the expression below, not a complete specification of the source problem's admissible class. An isolated singularity here refers to the source's intended singular point \(0\), with no other singularities sufficiently close to it; the relevant notion of singularity and regularity is not supplied.

The requested question is whether every admissible extremal solution in the intended, but unspecified, PDE setting satisfies
\[
\limsup_{\substack{x\to 0\\x\in\Omega\setminus\{0\}}}
|x|^2 f'\bigl(u^\star(x)\bigr)<+\infty.
\]
Equivalently, the requested upper bound is the existence of constants \(C\in[0,\infty)\) and \(r_0>0\), possibly depending on the solution and its data, such that
\[
\forall x\in\Omega,\qquad
0<|x|<r_0\ \Longrightarrow\
f'\bigl(u^\star(x)\bigr)\leq \frac{C}{|x|^2}.
\]
This is an ordinary pointwise limit superior, so a specified representative or sufficient regularity away from the singularity is required. The source does not request a uniform constant over all nonlinearities, domains, or solutions. Prove this assertion under the intended hypotheses, or disprove it by an admissible example; the missing hypotheses must first be recovered before either task has a fully determined meaning.

The supplied context attributes to Brezis a question about the significance of dimension \(10\) and smoothness of \(u^\star\) in low dimensions for every \(f\) and \(\Omega\) in the intended setting. It attributes to Cabr\'e, Figalli, Serra, and Ros-Oton a result described as boundedness of weak stable solutions for \(n\leq 9\), together with
\[
\nabla u\in L^{2+\gamma}_{\mathrm{loc}}(\Omega),\qquad \gamma=\gamma(n)>0.
\]
These are supplied contextual claims whose PDE hypotheses and qualifications are absent from the excerpt. The passage associates the requested estimate with Brezis' Open Question 4 and states that stability gives an averaged version and a ``lim-inf'' version. Its extracted averaged formula has the form \(\int f'(u)\leq Cr^{n-2}\), but lacks an integration region; the precise averaged and limit-inferior statements therefore cannot be reconstructed from this record.
\subsection{Short English statement}
The question asks whether, near an isolated singularity at \(0\), the quantity \(f'(u^\star(x))\) for an extremal solution has a pointwise upper bound of order \(|x|^{-2}\). The equation, admissible data, boundary conditions, and meaning of extremality are missing, so the displayed statement is a conditional formulation requiring human review. The supplied open label and literature assessment are dated curation, not independent verification.
\subsection{Sources}
\begin{itemize}
\item Local registry: \texttt{data/unsolvedmath/problems.json}, supplied record with permanent UnsolvedMath ID \(30005996\), exact source problem code \texttt{catalogue entry}, title ``Pointwise Quadratic Bounds for Extremal Solutions.'' The record labels the problem open and its statement \texttt{corrected\_verified}, and reports that reconstruction was checked against the official report. These are supplied metadata, not verification performed for this draft. No individual proposer is supplied.
\item Supplied report citation: \emph{Calculus of Variations (2025)}, Oberwolfach Reports, DOI \texttt{10.4171/owr/2024/37}, pp.\ 2113--2178. \url{https://doi.org/10.4171/owr/2024/37}. The supplied title uses 2025, while the report identifier and proposed-year metadata use 2024; both are retained without reconciliation.
\item Supplied literature source: ``Oberwolfach Report 37/2024,'' \url{https://ems.press/content/serial-article-files/50045}.
\item Supplied literature source: ``Figalli--Franceschini, 2026,'' \url{https://arxiv.org/abs/2606.21546}. No paper title or detailed theorem statement is supplied.
\item Literature assessment dated exactly \texttt{2026-08-22}: the supplied curation labels the problem open, says that low-dimensional regularity is proved, and says that pointwise finiteness of the requested limit superior remains unresolved. It describes recent work as giving averaged bounds and a limit-superior lower-bound mechanism, rather than the requested upper bound. This is supplied dated curation derived from the OWR material; its time-sensitive status and cited mathematical claims have not been independently verified for this draft.
\item Attribution within the supplied excerpt: Brezis is cited as reference [1], including Open Problem 1 and Open Question 4; Cabr\'e, Figalli, Serra, and Ros-Oton are cited as reference [3]. Earlier partial results are cited as [12, 7, 4, 6, 14, 8, 5]. The bibliography identifying these numbered references is not included in the supplied record.
\end{itemize}
\end{document}
